\documentclass[11pt]{article}

\usepackage{amsmath}
\usepackage{amssymb}
\usepackage{mathtools}
\usepackage{siunitx}
\usepackage[backend=biber,style=numeric]{biblatex}
\usepackage[colorlinks=true,linkcolor=blue,citecolor=blue,urlcolor=blue]{hyperref}

\addbibresource{references.bib}

\title{Provider Equation Documentation Mirror}
\author{ePC-SAFT Provider}
\date{2026-07-24}

\begin{document}
\maketitle

\section{Scope and authority}

This document mirrors the current Provider implementation. It is not a runtime
input and does not override \texttt{docs/science/full-provider-eos.md}. The
legacy catalogue at commit
\texttt{7920e56dbb0fc1fb418bae44e67ef50f334b75cd}, Git blob
\texttt{2e970476258558de05391d9d77f31d1736d78de8}, is provenance and structural
guidance only. CppAD is the sole nonlinear derivative authority.

\section{Symbols, bases, and units}

The residual molar Helmholtz energy is
$a^r=A^r/(nRT)$. Molar density $\rho_m$ is in
\si{\mole\per\meter\cubed}; segment diameters and size parameters are in
\si{\angstrom}; $T$ and $\epsilon_i/k_B$ are in kelvin. The symbol
$\widehat{\rho}$ denotes the corresponding number density in
\si{\angstrom\tothe{-3}}.

% EqID: angstrom_number_density
% Source: \cite{Gross2001}, unit conversion used by Appendix A
% Status: Exact derivation
% Description: Convert molar density to molecular number density on the angstrom basis.
% Change note: Uses the current Provider Avogadro constant and angstrom conversion.
\begin{equation}
\widehat{\rho}=\rho_m\,N_A\,10^{-30}.
\label{eq:angstrom_number_density}
\end{equation}

\section{Parameter setup}

% EqID: segment_sigma_correlation
% Source: Source-receipted component correlation represented by the Provider parameter schema
% Status: Bounded empirical relation
% Description: Evaluate a component segment-size correlation from its constant and exponential terms.
% Change note: Mirrors the general sparse correlation implemented by segment_sigma.
\begin{equation}
\sigma_i(T)=\sigma_{i,0}+\sum_{q\in\mathcal{Q}_i} A_q\exp(b_qT).
\label{eq:segment_sigma_correlation}
\end{equation}

% EqID: neutral_segment_diameter
% Source: \cite{Gross2001}, Eq.~(A.6)
% Status: Source equation
% Description: Temperature-dependent hard-chain diameter for neutral components.
% Change note: Reconciled to the current neutral diameter branch.
\begin{equation}
d_i(T)=\sigma_i(T)\left[1-0.12\exp\left(-\frac{3(\epsilon_i/k_B)}{T}\right)\right],
\qquad z_i=0.
\label{eq:neutral_segment_diameter}
\end{equation}

% EqID: figiel_ionic_segment_diameter
% Source: \cite{Figiel2025}, ionic hard-chain diameter convention
% Status: Implementation convention
% Description: Figiel-family ionic hard-chain diameter.
% Change note: Preserves the current 0.88 sigma routing for Figiel electrolyte models.
\begin{equation}
d_i=0.88\,\sigma_i,\qquad z_i\ne0,\quad
\mathcal{M}\ne\mathcal{M}_{\mathrm{Ascani}}.
\label{eq:figiel_ionic_segment_diameter}
\end{equation}

% EqID: ascani_ionic_segment_diameter
% Source: \cite{AscaniHeld2021}, ion diameter convention
% Status: Source equation
% Description: Ascani original-Born ionic hard-chain diameter.
% Change note: Preserves the source-routed full-sigma diameter.
\begin{equation}
d_i=\sigma_i,\qquad z_i\ne0,\quad
\mathcal{M}=\mathcal{M}_{\mathrm{Ascani}}.
\label{eq:ascani_ionic_segment_diameter}
\end{equation}

% EqID: cross_segment_diameter
% Source: \cite{Gross2001}, Eq.~(A.7), with source-receipted l_ij extension
% Status: Source equation
% Description: Unlike-pair segment diameter including the distinct l_ij family.
% Change note: Makes the current l_ij semantics explicit.
\begin{equation}
\sigma_{ij}=\frac{\sigma_i+\sigma_j}{2}(1-l_{ij}).
\label{eq:cross_segment_diameter}
\end{equation}

% EqID: cross_dispersion_energy
% Source: \cite{Gross2001}, Eq.~(A.7)
% Status: Source equation
% Description: Unlike-pair dispersion energy for an applicable pair.
% Change note: Keeps k_ij semantically distinct from l_ij and association data.
\begin{equation}
\frac{\epsilon_{ij}}{k_B}
=\sqrt{\frac{\epsilon_i}{k_B}\frac{\epsilon_j}{k_B}}\,(1-k_{ij}),
\qquad z_iz_j\le0.
\label{eq:cross_dispersion_energy}
\end{equation}

% EqID: same_sign_ionic_dispersion
% Source: \cite{Figiel2025}, same-sign ion convention
% Status: Implementation convention
% Description: Suppress dispersion between equally charged ions.
% Change note: Mirrors the charged-pair branch in residual_helmholtz.
\begin{equation}
\epsilon_{ij}=0,\qquad z_iz_j>0.
\label{eq:same_sign_ionic_dispersion}
\end{equation}

\section{Density and packing}

% EqID: mean_segment_count
% Source: \cite{Gross2001}, Appendix A mixture definition
% Status: Source equation
% Description: Composition-averaged segment count.
% Change note: Reconciled to the current residual Helmholtz owner.
\begin{equation}
\overline{m}=\sum_i x_i m_i.
\label{eq:mean_segment_count}
\end{equation}

% EqID: packing_moments
% Source: \cite{Gross2001}, Eq.~(A.8)
% Status: Source equation
% Description: Segment-density moments used by hard-chain and association terms.
% Change note: Uses the current angstrom number-density basis.
\begin{equation}
\zeta_n=\frac{\pi}{6}\widehat{\rho}\sum_i x_i m_i d_i^n,
\qquad n=0,1,2,3.
\label{eq:packing_moments}
\end{equation}

% EqID: packing_fraction
% Source: \cite{Gross2001}, Eq.~(A.8) at n=3
% Status: Source equation
% Description: Packing fraction from molar density and segment geometry.
% Change note: Shared by Helmholtz evaluation and Provider domain owners.
\begin{equation}
\eta=\zeta_3
=\frac{\pi}{6}\rho_mN_A10^{-30}\sum_i x_i m_i d_i^3.
\label{eq:packing_fraction}
\end{equation}

% EqID: close_packing_fraction
% Source: Geometric close packing of equal spheres
% Status: Exact derivation
% Description: Provider hard upper packing-domain boundary.
% Change note: Mirrors close_packing_fraction without introducing solver policy.
\begin{equation}
\eta_{\mathrm{cp}}=\frac{\pi}{3\sqrt{2}}.
\label{eq:close_packing_fraction}
\end{equation}

% EqID: packing_fraction_coefficient
% Source: Exact specialization of Eq.~\eqref{eq:packing_fraction}
% Status: Exact derivation
% Description: Composition and temperature coefficient satisfying eta=rho_m c_eta.
% Change note: Mirrors the value-only volume-domain helper.
\begin{equation}
c_\eta(T,\boldsymbol{x})=
\frac{\pi}{6}N_A10^{-30}\sum_i x_i m_i d_i(T)^3,
\qquad \eta=\rho_m c_\eta.
\label{eq:packing_fraction_coefficient}
\end{equation}

% EqID: phase_packing_fraction
% Source: Exact extensive-coordinate form of Eq.~\eqref{eq:packing_fraction}
% Status: Exact derivation
% Description: Packing fraction in explicit component amounts and extensive volume.
% Change note: Mirrors the native phase-domain derivative owner.
\begin{equation}
\eta(T,\boldsymbol{n},V)=
\frac{\pi}{6}\frac{N_A10^{-30}}{V}\sum_i n_i m_i d_i(T)^3.
\label{eq:phase_packing_fraction}
\end{equation}

\section{Contribution intermediates}

\subsection{Hard chain}

% EqID: hard_sphere_helmholtz
% Source: \cite{Gross2001}, Eqs.~(A.8)--(A.9)
% Status: Source equation
% Description: Boublik-Mansoori-Carnahan-Starling-Leland hard-sphere Helmholtz term.
% Change note: Mirrors the current zeta-based expression.
\begin{equation}
a^{hs}=\frac{1}{\zeta_0}\left[
\frac{3\zeta_1\zeta_2}{1-\zeta_3}
+\frac{\zeta_2^3}{\zeta_3(1-\zeta_3)^2}
+\left(\frac{\zeta_2^3}{\zeta_3^2}-\zeta_0\right)
\ln(1-\zeta_3)\right].
\label{eq:hard_sphere_helmholtz}
\end{equation}

% EqID: hard_sphere_contact
% Source: \cite{Gross2001}, Eq.~(A.17)
% Status: Source equation
% Description: Hard-sphere radial distribution function at segment contact.
% Change note: Uses the current pure-component contact diameter d_i/2.
\begin{equation}
g_{ii}^{hs}=
\frac{1}{1-\eta}
+\frac{3(d_i/2)\zeta_2}{(1-\eta)^2}
+\frac{2(d_i/2)^2\zeta_2^2}{(1-\eta)^3}.
\label{eq:hard_sphere_contact}
\end{equation}

% EqID: hard_chain_helmholtz
% Source: \cite{Gross2001}, Eq.~(A.16)
% Status: Source equation
% Description: PC-SAFT hard-chain residual Helmholtz contribution.
% Change note: Mirrors the site-independent chain term.
\begin{equation}
a^{hc}=\overline{m}\,a^{hs}
-\sum_i x_i(m_i-1)\ln g_{ii}^{hs}.
\label{eq:hard_chain_helmholtz}
\end{equation}

\subsection{Dispersion}

% EqID: dispersion_chain_coefficients
% Source: \cite{Gross2001}, Eqs.~(A.13)--(A.14)
% Status: Source equation
% Description: Chain-length interpolation of universal dispersion coefficients.
% Change note: Mirrors the current A0/A1/A2 and B0/B1/B2 arrays.
\begin{equation}
\begin{aligned}
a_k(\overline m)&=a_{0k}
+\frac{\overline m-1}{\overline m}a_{1k}
+\frac{(\overline m-1)(\overline m-2)}{\overline m^2}a_{2k},\\
b_k(\overline m)&=b_{0k}
+\frac{\overline m-1}{\overline m}b_{1k}
+\frac{(\overline m-1)(\overline m-2)}{\overline m^2}b_{2k}.
\end{aligned}
\label{eq:dispersion_chain_coefficients}
\end{equation}

% EqID: dispersion_polynomial
% Source: \cite{Gross2001}, polynomial form in Eqs.~(A.11)--(A.12)
% Status: Exact derivation
% Description: Horner-evaluated sixth-order dispersion polynomial.
% Change note: Documents the real arithmetic performed by polynomial.
\begin{equation}
\mathcal{P}(\boldsymbol{c},\eta)=\sum_{k=0}^{6}c_k\eta^k.
\label{eq:dispersion_polynomial}
\end{equation}

% EqID: dispersion_integral_one
% Source: \cite{Gross2001}, Eq.~(A.11)
% Status: Source equation
% Description: First PC-SAFT dispersion integral correlation.
% Change note: Evaluated through the shared polynomial owner.
\begin{equation}
I_1(\eta,\overline m)=\sum_{k=0}^{6}a_k(\overline m)\eta^k.
\label{eq:dispersion_integral_one}
\end{equation}

% EqID: dispersion_integral_two
% Source: \cite{Gross2001}, Eq.~(A.12)
% Status: Source equation
% Description: Second PC-SAFT dispersion integral correlation.
% Change note: Evaluated through the shared polynomial owner.
\begin{equation}
I_2(\eta,\overline m)=\sum_{k=0}^{6}b_k(\overline m)\eta^k.
\label{eq:dispersion_integral_two}
\end{equation}

% EqID: dispersion_compressibility_factor
% Source: \cite{Gross2001}, Eq.~(A.15)
% Status: Source equation
% Description: Chain compressibility correction multiplying the second-order term.
% Change note: Mirrors the current rational expression.
\begin{equation}
C_1=\left[
1+\overline m\frac{8\eta-2\eta^2}{(1-\eta)^4}
+(1-\overline m)
\frac{20\eta-27\eta^2+12\eta^3-2\eta^4}
{[(1-\eta)(2-\eta)]^2}
\right]^{-1}.
\label{eq:dispersion_compressibility_factor}
\end{equation}

% EqID: dispersion_first_moment
% Source: \cite{Gross2001}, Eq.~(A.10) mixture sum
% Status: Source equation
% Description: First unlike-pair dispersion mixture moment.
% Change note: Includes current segment counts and cross diameter.
\begin{equation}
M_1=\sum_i\sum_j x_ix_jm_im_j\sigma_{ij}^3
\frac{\epsilon_{ij}/k_B}{T}.
\label{eq:dispersion_first_moment}
\end{equation}

% EqID: dispersion_second_moment
% Source: \cite{Gross2001}, Eq.~(A.10) mixture sum
% Status: Source equation
% Description: Second unlike-pair dispersion mixture moment.
% Change note: Includes the squared reduced pair energy.
\begin{equation}
M_2=\sum_i\sum_j x_ix_jm_im_j\sigma_{ij}^3
\left(\frac{\epsilon_{ij}/k_B}{T}\right)^2.
\label{eq:dispersion_second_moment}
\end{equation}

% EqID: dispersion_helmholtz
% Source: \cite{Gross2001}, Eq.~(A.10)
% Status: Source equation
% Description: PC-SAFT first- and second-order dispersion Helmholtz contribution.
% Change note: Reconciled to the current angstrom number-density basis.
\begin{equation}
a^{disp}=-2\pi\widehat{\rho}I_1M_1
-\pi\widehat{\rho}\,\overline m\,C_1I_2M_2.
\label{eq:dispersion_helmholtz}
\end{equation}

\subsection{Association}

% EqID: association_contact
% Source: \cite{Gross2001}, hard-sphere contact rule used by association
% Status: Source equation
% Description: Unlike-site hard-sphere contact function.
% Change note: Mirrors the reduced-diameter association contact expression.
\begin{equation}
g_{ij}^{hs}=
\frac{1}{1-\eta}
+\frac{3d_id_j}{d_i+d_j}\frac{\zeta_2}{(1-\eta)^2}
+2\left(\frac{d_id_j}{d_i+d_j}\right)^2
\frac{\zeta_2^2}{(1-\eta)^3}.
\label{eq:association_contact}
\end{equation}

% EqID: association_strength
% Source: \cite{Gross2002}, cross-association rule; \cite{Chapman1990}, association strength
% Status: Source equation
% Description: Site-pair association strength using explicit source-resolved cross records.
% Change note: Keeps association energy and volume distinct from k_ij and l_ij.
\begin{equation}
\Delta_{AB}^{ij}=g_{ij}^{hs}
\left[\exp\left(\frac{\epsilon_{AB}^{ij}/k_B}{T}\right)-1\right]
\sigma_{ij}^3\kappa_{AB}^{ij}.
\label{eq:association_strength}
\end{equation}

% EqID: association_mass_action
% Source: \cite{Chapman1990}, Eq.~(19)
% Status: Source equation
% Description: Association site-fraction mass-action equation with explicit multiplicity.
% Change note: Written in the zero-residual form used by the Provider.
\begin{equation}
F_A=X_A\left[
1+\widehat{\rho}\sum_B x_{i(B)}\,\nu_B\,X_B\Delta_{AB}
\right]-1=0.
\label{eq:association_mass_action}
\end{equation}

% EqID: association_helmholtz
% Source: \cite{Chapman1990}, Eq.~(21)
% Status: Source equation
% Description: Association residual Helmholtz contribution.
% Change note: Uses explicit site multiplicity from the resolved model.
\begin{equation}
a^{assoc}=\sum_A x_{i(A)}\nu_A
\left(\ln X_A-\frac{X_A}{2}+\frac{1}{2}\right).
\label{eq:association_helmholtz}
\end{equation}

\section{Residual Helmholtz energy}

% EqID: residual_helmholtz
% Source: \cite{Gross2001}, PC-SAFT residual decomposition, extended by the admitted electrolyte terms
% Status: Implementation convention
% Description: Sum of every active residual Helmholtz contribution.
% Change note: Mirrors the single Provider Helmholtz owner and its contribution result.
\begin{equation}
a^r=a^{hc}+a^{disp}+a^{assoc}+a^{DH}+a^{Born}.
\label{eq:residual_helmholtz}
\end{equation}

\printbibliography

\end{document}
